\section{Speech: 5-minute version (9 slides)}

Use the deck \texttt{ISASF2026\_Paper132\_Presentation\_5min.pptx}. Speak slowly: the script is about
4 minutes 15 seconds, which leaves time for changing slides. If you run short of time, skip the sentences marked
\textit{(optional)}.

\slide{1}{Title}{15 s}
\say{Good morning. I am [name]. Our paper asks one question: can a small vision-language model, fine-tuned on a
single 8-gigabyte GPU, turn a screenshot of a user interface into React or Vue code?}

\slide{2}{Introduction}{40 s}
\say{Screenshot-to-code research mostly produces HTML, but real frontends are built from React and Vue components,
and large commercial models cost money per token and see your unreleased designs.
No dataset pairs screenshots with validated React and Vue code. And small models struggle here: used zero-shot,
the 2B model imports UI libraries in 437 of its 496 React outputs, even though we asked for Tailwind only.}

\slide{3}{Methodology}{50 s}
\say{We did three things. First, we extended WebSight: a code model generated React and Vue versions of every page,
each checked by a real parser; 15,588 samples have both targets valid.
Second, we fine-tuned two models with QLoRA on one 8-gigabyte GPU: Qwen2-VL-2B, and for comparison Qwen2.5-VL-3B.
Third, we evaluated the code, and we rendered every output in a browser to compare it with the original screenshot,
with bootstrap confidence intervals.}

\slide{4}{Results: WebSight}{55 s}
\say{On 1,488 test samples, fine-tuning the 2B model raises BLEU-4 from 52 to 81 and CodeBLEU from 0.52 to 0.74.
Against the larger zero-shot 3B model, our renders are significantly closer on LPIPS and PSNR.
On SSIM the two are statistically tied; SSIM is near its ceiling of 0.722.
The fine-tuned 3B improves a lot too, but does not beat the 2B. \textit{(optional:) It had only one epoch instead of three.}}

\slide{5}{Efficiency}{35 s}
\say{Our 2B model needs 2.18 gigabytes of GPU memory, a third of the zero-shot 3B model, and it is about twice as fast.
But to be clear: 4-bit QLoRA saves memory, not time. It is 1.75 times slower than the 16-bit base model.}

\slide{6}{Real web pages (Design2Code)}{45 s}
\say{Reviewers asked about real websites, so we tested 100 real pages from Design2Code.
Fine-tuning still helps strongly: SSIM rises by 0.21 over the base model.
But here our model is only tied with the 3B models, and it writes about 600 tokens per page, while real HTML has a
median of nearly 15,000. It learned the length of the short synthetic training pages.}

\slide{7}{Conclusion}{30 s}
\say{To conclude: we built more than 15,000 parser-validated React and Vue targets, and showed that a 2B model
fine-tuned on one consumer GPU can match a larger fine-tuned model and, on WebSight-style pages, beat a larger zero-shot
one, with designs kept on-premises.
\textit{(optional:) Next, we want to train on longer real pages and compare QLoRA with LoRA and DoRA.}}

\slide{8--9}{References and Thank You}{5 s}
\say{Thank you. I am happy to take your questions.}

\medskip
\noindent\textbf{Timing check:} 15 + 40 + 50 + 55 + 35 + 45 + 30 + 5 seconds = 4 minutes 35 seconds.
The longer 14-slide deck (\texttt{ISASF2026\_Paper132\_Presentation.pptx}) stays as a backup; the slides removed from
the 5-minute version (abstract, literature review, dataset, discussion, limitations) are useful if someone asks.

